% Source: https://artofproblemsolving.com/community/c3h2448723_1729_posts_also_problems
% Stable identity: aops-topic-2448723-q01
In a parallelogram the opposite sides are parallel and equal: $AB\parallel CD$ and $AD=BC=10$. Since $E$ lies on line $CD$, segment $DE$ is parallel to $AB$.

\begin{asy}
usepackage("amsmath");
size(7cm);
defaultpen(fontsize(10pt));
pair A=(0,0),B=(16,0),C=B+10*dir(60),D=A+(C-B),E1=D+(-4,0);
pair F=extension(B,E1,A,D);
fill(A--B--F--cycle,rgb(0.85,0.92,1)); fill(D--E1--F--cycle,rgb(1,0.88,0.85));
draw(A--B--C--D--cycle,linewidth(0.9)); draw(D--E1,dashed); draw(B--E1,red);
draw(arc(B,1.6,60,180)); label("$120^\circ$",B+2.6*dir(120),fontsize(8pt));
dot(F); dot(E1);
label("$A$",A,SW);label("$B$",B,SE);label("$C$",C,NE);label("$D$",D,N);label("$E$",E1,NW);label("$F$",F,dir(160));
label("$16$",(A+B)/2,S);label("$10$",(B+C)/2,E);label("$4$",(D+E1)/2,N);
\end{asy}

In triangles $FAB$ and $FDE$, the angles at $F$ are vertical angles, and $\angle FAB=\angle FDE$ because they are alternate angles for the parallel lines $AB$ and $DE$. So $\triangle FAB\sim\triangle FDE$ with ratio $\dfrac{AB}{DE}=\dfrac{16}{4}=4$, hence $AF=4\,FD$. Since $F$ lies on segment $AD$,
\[AF+FD=AD=10\implies5\,FD=10\implies FD=2.\]
So $FD$ is exactly $2$, answer $\boxed{\textbf{(B) }2}$. (The angle $120^\circ$ is not needed.)
